% ------------------------------------------------------------
% A6.x Substrate Uplift Module C — Unistochastic Emergence (v0.1)
% Author: Borealis S. Hedling
% Altitude: A6.x (Substrate Uplift Layer)
% Purpose: Define the unistochastic emergence substrate required
%          for spectral continuity, modal lift, and collapse-
%          compatible probabilistic geometry at A6 altitude.
% ------------------------------------------------------------

\documentclass[11pt]{article}
\usepackage{amsmath, amssymb, amsthm}

\title{A6.x Substrate Uplift Module C: Unistochastic Emergence}
\author{Borealis S. Hedling}
\date{October 2026}

\begin{document}
\maketitle

\section*{Overview}

This module defines the unistochastic emergence substrate required for A6
altitude geometry. Unistochastic operators provide the probabilistic continuity
and modal-preserving transition structure required by the A6 Constructive Fixes
and the A7 Composite Substrate Operator Algebra.

% ------------------------------------------------------------
\section{Unistochastic Operator Definition}

Let $U$ be a unitary matrix.

\textbf{Definition (Unistochastic Operator).}
\[
B_{ij} = |U_{ij}|^2.
\]

The matrix $B$ is unistochastic.

\textbf{Theorem (Stochasticity).}
\[
\sum_j B_{ij} = 1, \quad \sum_i B_{ij} = 1.
\]

Thus $B$ defines a probability-preserving operator.

% ------------------------------------------------------------
\section{Unistochastic Continuity}

Let $U(p)$ vary smoothly over the manifold.

\textbf{Theorem (Continuity).}
If $U(p)$ is smooth, then:
\[
B(p) \text{ continuous in } p.
\]

This enforces Fix 5 (spectral continuity).

% ------------------------------------------------------------
\section{Spectral Compatibility}

Define the spectral operator:
\[
\mathcal{S}_{\mathrm{uni}}(p) = \mathrm{Spec}(B(p)).
\]

\textbf{Theorem (Spectral Mismatch Bound).}
\[
\|\mathcal{S}_{\mathrm{uni}}(p) - \mathcal{S}_{\mathrm{uni}}(q)\| < \epsilon.
\]

This enforces Fix 6 (spectral mismatch bound).

% ------------------------------------------------------------
\section{Modal Lift Compatibility}

Let $\mu$ be the modal map.

\textbf{Definition (Unistochastic Modal Lift).}
\[
\mu(p) \leadsto \mu(q) \quad \text{if} \quad B_{pq} > 0.
\]

\textbf{Theorem (Lift Preservation).}
If $B_{pq} > 0$, then:
\[
\mu(p) \leadsto \mu(q).
\]

This enforces Fix 10 (modal lift preservation).

% ------------------------------------------------------------
\section{Collapse-Surface Compatibility}

Let $\Sigma$ be the collapse surface.

\textbf{Definition (Collapse-Compatible Unistochasticity).}
\[
p \in \Sigma \Rightarrow B_{pp} = 1.
\]

\textbf{Theorem (Local Detectability).}
If $B_{pp} = 1$, then:
\[
p \in \Sigma \Rightarrow \exists \epsilon > 0 :
B_\epsilon(p) \cap \Sigma \neq \emptyset.
\]

This enforces Fix 8 (collapse locality).

% ------------------------------------------------------------
\section{Stability Envelope}

Define:
\[
E_{\mathrm{uni}}(p) = \|B(p)\|.
\]

\textbf{Theorem (Monotonicity).}
If $E(p) \le E(q)$, then:
\[
E_{\mathrm{uni}}(p) \le E_{\mathrm{uni}}(q).
\]

This enforces Fix 9 (stability envelope monotonicity).

% ------------------------------------------------------------
\section{Conclusion}

The unistochastic emergence substrate defined here provides the probabilistic
continuity, modal lift preservation, spectral stability, and collapse
compatibility required for A6 Constructive Fixes and A7 Composite Substrate
Operator Algebra. It is the third foundational substrate module in the A6.x
uplift sequence.

\end{document}
